\documentclass[11pt,a4paper]{article}
\usepackage[utf8]{inputenc}
\usepackage[T1]{fontenc}
\usepackage{amsfonts}
\usepackage[french]{babel}
\frenchbsetup{StandardLists=true}
\usepackage{enumitem}
\usepackage{amssymb}
\usepackage{amsmath}
\usepackage[left=2cm,right=2cm,top=1.5cm,bottom=1.5cm]{geometry}
\usepackage{multirow}
\usepackage[dvipsnames]{xcolor}

\usepackage[T1]{fontenc}
\usepackage{fouriernc}
\usepackage{graphicx}

\usepackage{setspace}
\singlespacing

\usepackage{tcolorbox}
\usepackage{multicol}
\usepackage{caption}
\usepackage{fancyhdr}
\pagestyle{fancy}
\renewcommand\headrulewidth{1pt}
\fancyhead[L]{Lycée Denis Diderot }
\fancyhead[R]{Classe de 3PM}
\setcounter{page}{1}\fancyfoot[C]{page \thepage}
\fancyfoot[R]{Année 2025-2026}
\fancyfoot[L]{Alexis RUIZ}
\usepackage{multirow,tabularx,booktabs}

\usepackage{awesomebox}
\setlength{\aweboxvskip}{0mm}
\usepackage{pifont}
\usepackage{chemmacros}
\chemsetup{modules=all}
\setlength{\parindent}{0pt}

\renewcommand{\arraystretch}{1.6}

\frenchbsetup{StandardLists=true}

\begin{document}

\begin{center}
 \LARGE{Chapitre 5 - Le pH -- Exercices}
\end{center}

\normalsize

\hrule\
\\[0,2cm]

\textbf{Exercice 1 -- Acide, neutre ou basique ?} \\

On mesure le pH de différentes solutions. Compléter le tableau suivant :

\vspace{3mm}

\begin{center}
\begin{tabular}{|l|c|c|}
\hline
\textbf{Solution} & \textbf{pH} & \textbf{Caractère (acide / neutre / basique)} \\
\hline
Eau de piscine & 7,2 & \\
\hline
Soda au cola & 2,5 & \\
\hline
Liquide vaisselle & 9 & \\
\hline
Eau distillée & 7 & \\
\hline
Jus de pomme & 3,5 & \\
\hline
Déboucheur & 13,5 & \\
\hline
\end{tabular}
\end{center}

\textbf{Exercice 2 -- Classer des produits} \\[1mm]

On dispose de cinq produits du quotidien avec leur pH :

\vspace{1mm}

\noindent%
\begin{minipage}{0.32\linewidth}%
{\renewcommand{\arraystretch}{1.4}
\begin{tabular}{|l|c|}
\hline
\textbf{Produit} & \textbf{pH} \\
\hline
Eau minérale Evian & 7 \\
\hline
Coca-Cola & 2,5 \\
\hline
Savon de Marseille & 10 \\
\hline
Jus de citron pressé & 2 \\
\hline
Lait demi-écrémé & 6,5 \\
\hline
\end{tabular}}
\end{minipage}%
\hfill%
\begin{minipage}{0.63\linewidth}%
\begin{enumerate}[label=\textbf{\arabic*.}]
\item Classer ces produits du plus acide au plus basique.

\vspace{5mm}
\dotfill

\item Quel produit est neutre ? Justifier.

\vspace{5mm}
\dotfill

\item Y en a-t-il de corrosifs et dangereux ?

\vspace{5mm}
\dotfill
\end{enumerate}
\end{minipage}

\vspace{3mm}

\textbf{Exercice 3 -- Dilution et cuisine} \\[1mm]

Le jus de tomate a un pH d'environ 4. On le dilue en ajoutant de l'eau.

\begin{enumerate}[label=\textbf{\arabic*.}]
\item Le jus de tomate est-il acide, neutre ou basique ?

\vspace{5mm}
\dotfill

\item Comment va évoluer son pH ? Vers quelle valeur tend-il ?

\vspace{5mm}
\dotfill

\item Son pH peut-il dépasser 7 par dilution ? Justifier.

\vspace{5mm}
\dotfill

\item En cuisine, on ajoute souvent du sucre dans la sauce tomate. Expliquer pourquoi.

\vspace{5mm}
\dotfill

\vspace{5mm}

\dotfill
\end{enumerate}

\newpage

\textbf{Exercice 4 -- Les ions} \\

\begin{enumerate}[label=\textbf{\arabic*.}]
\item Quel est l'ion responsable du caractère acide d'une solution ? Donner son nom et sa formule.

\vspace{8mm}
\dotfill

\item Quel est l'ion responsable du caractère basique d'une solution ? Donner son nom et sa formule.

\vspace{8mm}
\dotfill

\item Une solution a un pH = 4. Quel ion est majoritaire dans cette solution : \ch{H+} ou \ch{HO-} ? Justifier.

\vspace{8mm}
\dotfill

\vspace{8mm}
\dotfill
\end{enumerate}

\vspace{3mm}

\textbf{Exercice 5 -- Indicateurs colorés} \\

On verse quelques gouttes de bleu de bromothymol (BBT) dans trois solutions A, B et C. On observe les couleurs suivantes :

\vspace{3mm}

\begin{center}
\begin{tabular}{|c|c|c|}
\hline
\textbf{Solution A} & \textbf{Solution B} & \textbf{Solution C} \\
\hline
Jaune & Bleu & Vert \\
\hline
\end{tabular}
\end{center}

\vspace{3mm}

\begin{enumerate}[label=\textbf{\arabic*.}]
\item Pour chaque solution, indiquer si elle est acide, neutre ou basique.

\vspace{8mm}
\dotfill

\vspace{8mm}
\dotfill

\vspace{8mm}
\dotfill

\item La zone de virage du BBT est 6,0 -- 7,6. Que peut-on dire du pH de la solution A ?

\vspace{8mm}
\dotfill

\dotfill
\end{enumerate}


\end{document}
